\documentclass[12pt]{article}

% --------------------
% PACOTES
% --------------------
\usepackage[utf8]{inputenc}
\usepackage[T1]{fontenc}
\usepackage[brazil]{babel}

\usepackage{amsmath, amssymb}
\usepackage{graphicx}
\usepackage{geometry}
\usepackage{hyperref}
\usepackage{booktabs}

\geometry{margin=2.5cm}

% --------------------
% COMANDOS
% --------------------
\newcommand{\ema}{Estações Meteorológicas Automáticas}
\newcommand{\mqtt}{MQTT}
\newcommand{\api}{API}

% --------------------
% DOCUMENTO
% --------------------
\title{Relatório Técnico \\ Sistema de \ema}
\author{Henry Campos}
\date{\today}

\begin{document}

\maketitle
\tableofcontents
\newpage

% --------------------
\section{Introdução}

Este relatório apresenta o desenvolvimento do sistema de \ema\ para monitoramento climático urbano.

% --------------------
\section{Arquitetura do Sistema}

O sistema é composto por três camadas principais:

\begin{itemize}
    \item Coleta de dados (estações)
    \item Backend (processamento e armazenamento)
    \item Clientes (visualização)
\end{itemize}

% --------------------
\section{Implementação}

A comunicação foi realizada utilizando o protocolo \mqtt, com backend baseado em serviços web.

% --------------------
\section{Resultados}

Os dados coletados permitiram visualizar variações climáticas locais em tempo real.

% --------------------
\section{Conclusão}

O sistema demonstrou viabilidade para aplicação em monitoramento urbano.

% --------------------
\bibliographystyle{plain}
\bibliography{referencias}

\end{document}
